\documentclass[11pt]{article}

\usepackage[margin=1in]{geometry}
\usepackage{array}
\usepackage{longtable}
\usepackage{hyperref}

\title{OmegaSim Experiment Rundown and Design Implications}
\author{Prepared for Ben Goertzel}
\date{2026-06-27}

\begin{document}
\maketitle

\begin{abstract}
This report summarizes the OmegaSim experiments run so far, the results
obtained, and the lessons for designing new experiments aimed at eliciting
more complex collective dynamics. The central historical lesson is
conservative: many apparent lobe, synchrony, or attractor-like effects were
explained by queue load, service capacity, action opportunity, transfer timing,
or other accounting variables. The central forward-looking opportunity is to
move from load-pressure and linear-prediction scaffolds toward explicitly
nonlinear, delayed, bounded-resource dependencies among agents or hives. In
particular, logistic-shaped dependence of one agent's activity on another
agent's lagged activity is a plausible ingredient for structured recurrent
dynamics, provided it is embedded in a system with thresholds, delays,
imperfect prediction, resource costs, heterogeneity, and controls that separate
true residual dynamics from queueing arithmetic.
\end{abstract}

\section{Purpose}

This document is meant as a compact but detailed handoff to a strategic
reviewer or collaborator. It records what OmegaSim has already tested, which
claims survived controls, which claims failed, and what the existing evidence
implies for future logistic-dependence experiments.

The detailed source artifacts live primarily under:

\begin{verbatim}
docs/results/
docs/a5_anticipatory_predictive_control_preregistration.md
docs/structured_strange_attractor_research_note.md
\end{verbatim}

\section{Current Simulator Abstraction}

OmegaSim is an abstract deterministic simulator for OmegaHive-like and
Moltbook-like dynamics. It does not call real LLMs, Lean, Slack, browsers,
Atomspace, live task boards, or external services. Those systems are
represented numerically through queues, task classes, role policies, graph
edges, events, and metrics.

\subsection{A0/A1 Baseline}

The baseline system contains:

\begin{itemize}
\item 15 static agents.
\item Five roles, three agents per role: coordinator, researcher, architect,
implementer, and reviewer.
\item One NetworkX bus graph.
\item One in-memory task queue in the single-hive case.
\item Four baseline actions: idle, message, create\_task, and work\_task.
\item Deterministic YAML configuration and seeded runs.
\item Stable artifacts: config, manifest, metrics, events, and summary.
\end{itemize}

The baseline lobe labels are:

\begin{verbatim}
backlog_growth
execution
task_generation
coordination
low_activity
\end{verbatim}

These labels are derived from queue movement and dominant action counts. This
was useful for early trajectory summaries, but later became an interpretive
limitation: because the labels depend directly on queue and action fields, they
are easily confounded with task inflow, service capacity, and action
opportunity.

\subsection{A2 Attention Classes}

A2 added an opt-in attention policy over four task classes:

\begin{verbatim}
near_term_external
long_term_research
internal_improvement
housekeeping
\end{verbatim}

The policy can use quota balancing or a seeded random-available scheduler. Runs
report per-class queues, completions, queued age, attention shares,
value-weighted completions, capture pressure, and effort-normalized value
fields.

\subsection{A4 Multi-Hive Scaffold}

A4 introduced an opt-in multi-hive form with per-hive queues, events, metrics,
and coupling modes:

\begin{verbatim}
none
direct
delayed
shuffled
\end{verbatim}

Coupling events and cross-hive metrics track transfer attempts, completed
transfers, inbound/outbound transfer counts, aggregate queue balance,
queued-age divergence, and completion-fraction divergence or correlation.

\subsection{A5 Predictive-Control Scaffold}

A5 added a single-hive predictive-control scaffold. It keeps demand streams,
task volume, action opportunity, service capacity, and work budget matched
while varying prediction condition:

\begin{verbatim}
reactive
linear
nonlinear
oracle
linear-budget shuffled
nonlinear-budget shuffled
\end{verbatim}

Forecast skill is treated as a manipulation check, not as evidence of emergent
dynamics by itself.

\section{Experiment Timeline and Results}

\subsection{A0/A1: Deterministic Baseline Harness}

\paragraph{Question.}
Can the project produce deterministic OmegaHive-like run artifacts with stable
schemas, reproducible seeds, and enough metrics for later dynamics experiments?

\paragraph{Result.}
Yes. A0/A1 established the deterministic runner, 15-agent hive, task queue, bus
graph, event stream, metrics, summaries, and reproducibility tests. It did not
make a scientific claim about complex dynamics.

\paragraph{Lesson.}
The artifact discipline is worth preserving. However, baseline lobe labels are
too directly tied to queue and action variables to support a strange-attractor
claim without stronger controls or new observables.

\subsection{A2.1: Attention-Policy and Pressure Exploration}

\paragraph{Question.}
Do attention policies and task-creation pressure generate distinct trajectory
regimes or lobe-like structure?

\paragraph{Design.}
A2 compared quota-balanced attention, research-heavy attention, and
internal-improvement-heavy attention under normal, medium, high, and later
extreme task-creation pressure. Early blocks used seeds 1--3; a larger pressure
baseline used seeds 1--69.

\paragraph{Main result.}
The robust finding was pressure-amplified queue depth. In the extreme-pressure
baseline over seeds 1--69, the strongest global response was final queue depth
for the internal-improvement policy:

\begin{center}
\begin{tabular}{lrrrr}
condition means & normal & medium & extreme & high-minus-normal \\
\hline
final queue depth & 19.681159 & 45.956522 & 56.057971 & 36.376812
\end{tabular}
\end{center}

This is a real simulator effect, but it is primarily a pressure/load effect. It
does not by itself demonstrate emergent lobe grammar.

\paragraph{Secondary signals.}
Capture-pressure and value-yield fields showed policy-dependent structure. For
example, baseline peak long-term-research capture pressure declined from
0.282358 to 0.133015 across normal to extreme pressure. These signals remained
secondary and did not carry the main scientific claim.

\paragraph{Interpretation.}
Task-creation pressure is a strong stress knob. It changes backlog, queued age,
completion fraction, dwell, and transition summaries. Pressure alone mostly
produces queueing arithmetic rather than independent complex dynamics.

\subsection{A2.2: Scheduler Ablation}

\paragraph{Question.}
Is the pressure-amplified backlog response caused by the quota-balancing
scheduler?

\paragraph{Design.}
Quota balance was compared with seeded random-available selection. A pilot used
seeds 1--3. A preregistered holdout used seeds 70--99.

\paragraph{Holdout result.}
Mean final queue depth rose strongly under both schedulers:

\begin{center}
\begin{tabular}{lrrr}
scheduler & normal & medium & extreme \\
\hline
quota balance & 21.5 & 48.2 & 57.733333 \\
random available & 21.6 & 45.5 & 56.166667
\end{tabular}
\end{center}

The paired scheduler effect, random minus quota, was small compared with the
pressure effect:

\begin{center}
\begin{tabular}{lrr}
pressure level or contrast & mean delta & bootstrap interval \\
\hline
normal & 0.1 & [-2.233333, 2.066667] \\
medium & -2.7 & [-4.833333, -0.766667] \\
extreme & -1.566667 & [-4.3, 0.833333] \\
extreme-minus-normal & -1.666667 & [-5.1, 1.633333]
\end{tabular}
\end{center}

Load-normalized normal-to-extreme queue growth was nearly identical:

\begin{center}
\begin{tabular}{lrr}
scheduler & queue growth per added created task & bootstrap interval \\
\hline
quota balance & 1.209774 & [1.154681, 1.249447] \\
random available & 1.200961 & [1.145654, 1.259360]
\end{tabular}
\end{center}

\paragraph{Interpretation.}
The scheduler changes curve shape and secondary metrics, but it does not remove
the main pressure response. This shifted the interpretation away from a
scheduler-driven lobe mechanism and toward pressure plus finite service
capacity as the dominant explanation.

\subsection{A2.3: Exogenous Arrival Decoupling}

\paragraph{Question.}
Can backlog pressure be induced by external demand without increasing
agent-created tasks?

\paragraph{Design.}
The exogenous-arrival holdout used seeds 70--99, fixed task-creation pressure
at 1.0, quota balancing, baseline attention shares, and four exogenous arrival
rates.

\paragraph{Load result.}

\begin{center}
\begin{tabular}{lrrrr}
condition & agent-created & exogenous-created & total-created & queue per created \\
\hline
endogenous control & 41.433333 & 0.000000 & 41.433333 & 0.507770 \\
exogenous low & 40.966667 & 11.666667 & 52.633333 & 0.600270 \\
exogenous medium & 40.500000 & 23.266667 & 63.766667 & 0.682632 \\
exogenous high & 40.166667 & 38.100000 & 78.266667 & 0.749526
\end{tabular}
\end{center}

High exogenous arrivals raised total created tasks from 41.433333 to 78.266667
while agent-created tasks stayed near control.

\paragraph{Trajectory result.}
Queue-blind entropy declined only modestly at high exogenous demand:

\begin{center}
\begin{tabular}{lrr}
condition & queue-blind entropy & delta vs control \\
\hline
endogenous control & 2.353070 & 0.000000 \\
exogenous high & 2.254402 & -0.098668
\end{tabular}
\end{center}

\paragraph{Interpretation.}
Exogenous demand reproduces backlog pressure while decoupling the effect from
agent create-task pressure. This supports a load/accounting interpretation. The
residual queue-blind entropy signal is weak and descriptive, not a strong
lobe-grammar finding.

\subsection{A2.4: Demand versus Service Capacity}

\paragraph{Question.}
Does changing service capacity explain queue and trajectory effects that might
otherwise look like lobe dynamics?

\paragraph{Design.}
Seeds 70--99 crossed creation pressure 1.0, 1.8, and 2.2 with work service
capacity 0.7, 1.0, and 1.3.

\paragraph{Fixed-pressure service-capacity effects.}
Raising service capacity from 0.7 to 1.3 reduced backlog at every pressure:

\begin{center}
\begin{tabular}{lrrrr}
fixed pressure & queue delta & queue/created delta & completion frac delta & age delta \\
\hline
normal & -16.433334 & -0.315145 & 0.315145 & -1.601838 \\
high & -16.533334 & -0.168067 & 0.168067 & -0.760996 \\
extreme & -16.900000 & -0.151263 & 0.151263 & -0.633885
\end{tabular}
\end{center}

\paragraph{Fixed-service pressure effects.}
Raising creation pressure from 1.0 to 2.2 increased backlog at every service
level:

\begin{center}
\begin{tabular}{lrrrr}
fixed service & queue delta & queue/created delta & completion frac delta & age delta \\
\hline
low & 34.900000 & 0.205410 & -0.205410 & 0.769206 \\
baseline & 36.233333 & 0.298512 & -0.298512 & 1.165276 \\
high & 34.433334 & 0.369292 & -0.369292 & 1.737159
\end{tabular}
\end{center}

\paragraph{Queue-blind lobe robustness.}
Queue-blind labels excluded queue depth and queue delta and used action fields
such as tasks worked, tasks created, messages, and idle. Even then, pressure
shifted the system toward task generation and longer dwell:

\begin{center}
\begin{tabular}{lrrr}
fixed service & entropy delta & dwell max delta & task-generation share delta \\
\hline
low service & -0.556783 & 2.100000 & 0.472222 \\
baseline service & -0.506507 & 2.066667 & 0.491667 \\
high service & -0.176233 & 0.766667 & 0.472222
\end{tabular}
\end{center}

\paragraph{Interpretation.}
Service capacity explains much of the queue and lobe-like behavior. Under
pressure, trajectories lock into fewer transitions and longer dwell; extra
service partially offsets this. This is important, but still fits a
queue-flow/service explanation better than a free-standing attractor claim.

\subsection{A3: Queue-Flow and Service Synthesis}

\paragraph{Question.}
What is the cleanest single-hive interpretation after the A2 controls?

\paragraph{Result.}
A3 froze the single-hive interpretation as queue-flow/service accounting. The
robust claims are:

\begin{itemize}
\item Demand pressure increases total inflow, load-normalized backlog, and
queued-task age.
\item Higher service capacity reduces backlog and improves completion fraction.
\item Exogenous arrivals can induce backlog pressure without increasing
agent-created tasks.
\item Same-tick queue-balance correlations are artifact/accounting diagnostics,
not synchronization evidence.
\item Lagged service/completion and service/load-change correlations exist but
are modest and not seed-stable enough for mechanism claims.
\end{itemize}

The strongest lagged service/completion condition-mean shift was 0.168460, with
reported seed median 0.220475, bootstrap median interval [-0.063508, 0.447535],
and sign stability 0.666667. This remained exploratory.

\paragraph{Decision.}
A2/A3 residual lobe grammar was closed under the current single-hive simulator
and label scheme.

\subsection{A4: Two-Hive Coupling}

\paragraph{Question.}
When two hives are coupled through a shared transfer layer, do service/load
phase relations show coordination effects beyond queue-flow and accounting?

\paragraph{Design.}
The A4 holdout used seeds 100--129 and four preregistered modes: none, direct,
delayed, and shuffled. The analyzer consumed 30 paired seeds, 240 per-hive
endpoint rows, and 120 cross-hive endpoint rows.

\paragraph{Transfer-volume result.}
Direct, delayed, and shuffled all added full-transfer coupling versus none:

\begin{center}
\begin{tabular}{lrrr}
comparison & transfer-volume delta & bootstrap interval & sign stability \\
\hline
direct minus none & 865.333333 & [857.700000, 873.533333] & 1.000000 \\
delayed minus none & 865.900000 & [857.066667, 874.733333] & 1.000000 \\
shuffled minus none & 865.333333 & [857.200000, 873.733333] & 1.000000 \\
direct minus shuffled & 0.000000 & [0.000000, 0.000000] & 0.000000
\end{tabular}
\end{center}

Direct minus shuffled was zero because with two hives each source has only one
legal non-source target. Thus shuffled was a schema/source-opportunity check,
not a meaningful phase-randomization null.

\paragraph{Initial cross-hive endpoints.}
Delayed minus none showed strong completion-fraction correlations:

\begin{center}
\begin{tabular}{lrrr}
endpoint & mean delta & bootstrap interval & sign stability \\
\hline
completion-fraction corr lag 0 & 0.678330 & [0.494969, 0.857877] & 1.000000 \\
completion-fraction corr lag 2 & 0.619752 & [0.416347, 0.802583] & 1.000000 \\
load-backlog corr lag 0 & 0.325600 & [0.118204, 0.517629] & 0.999000
\end{tabular}
\end{center}

This was interesting but still needed null and accounting controls.

\paragraph{Temporal null.}
The delayed-coupling temporal null found that completion-fraction correlations
survived circular-shift controls:

\begin{center}
\begin{tabular}{lrrr}
endpoint & observed minus null & outside null CI & positive rate \\
\hline
completion-fraction corr lag 0 & 0.719486 & true & 0.966667 \\
completion-fraction corr lag 2 & 0.657053 & true & 0.900000 \\
load-backlog corr lag 0 & 0.333787 & false & 0.700000 \\
queued-age divergence final & -0.570895 & false & 0.433333
\end{tabular}
\end{center}

\paragraph{Accounting-control analyzer.}
The follow-up residualized completion-fraction trajectories by clock trend,
opportunity/load, action mix, and transfer timing, then reran circular-shift
nulls. Under combined non-tautological controls, the completion-fraction
correlations were no longer outside the null:

\begin{center}
\begin{tabular}{lrrrr}
endpoint & observed delta & null mean & null interval & outside CI \\
\hline
lag 0 & 0.156653 & -0.008884 & [-0.335280, 0.312277] & false \\
lag 2 & 0.140351 & -0.002389 & [-0.362347, 0.311764] & false
\end{tabular}
\end{center}

\paragraph{Decision.}
A4 was closed as queue-flow/service and action-opportunity accounting plus a
documented delayed synchronization diagnostic. It did not justify a three-hive
implementation, independent lobe grammar, or a strange-attractor claim.

\subsection{A5: Anticipatory Predictive Control}

\paragraph{Question.}
Can bounded prediction of future demand create richer residual collective
dynamics than pure reactivity or oracle-like smoothing?

\paragraph{Core hypothesis.}
Perfect prediction may smooth dynamics; zero prediction may leave the system
purely reactive. The interesting regime may be intermediate prediction budget,
where agents predict enough to find high-level recurring patterns but not
enough to eliminate errors. Those errors can create collective dynamics which
then become something new to predict.

\paragraph{Pilot, seeds 5--6.}
The pilot used matched task arrival, service, action, and demand-signal
settings. Mean final task creation was 78.5 tasks in every condition.

\begin{center}
\begin{tabular}{lrrrr}
condition & budget & forecast skill & completion frac & queue depth \\
\hline
reactive & 0.00 & 0.877957 & 0.386416 & 48.0 \\
linear & 0.35 & 0.911474 & 0.380463 & 48.5 \\
nonlinear & 0.65 & 0.961990 & 0.380463 & 48.5 \\
oracle & 1.00 & 1.000000 & 0.393265 & 47.5 \\
shuffled & 0.35 & 0.866049 & 0.407861 & 46.5
\end{tabular}
\end{center}

Intermediate predictors improved forecast skill, but primary residual-state
predictability contrasts did not survive paired label-permutation intervals.
The pilot failed closed for promotion.

\paragraph{Confirmatory run, seeds 7--16.}
The confirmatory run used 60 deterministic paired simulations across reactive,
linear, nonlinear, oracle, linear-budget shuffled, and nonlinear-budget
shuffled conditions. Mean final task creation was matched at 74.2 tasks in
every condition.

\begin{center}
\begin{tabular}{lrrrrr}
condition & budget & forecast skill & completion frac & queue depth & queued age \\
\hline
reactive & 0.00 & 0.877957 & 0.434071 & 42.7 & 5.725288 \\
linear & 0.35 & 0.911474 & 0.428319 & 43.1 & 5.774823 \\
nonlinear & 0.65 & 0.961990 & 0.420349 & 43.7 & 5.801917 \\
oracle & 1.00 & 1.000000 & 0.426118 & 43.3 & 5.783507 \\
shuffled & 0.35 & 0.866049 & 0.436788 & 42.6 & 5.681533 \\
nonlinear shuffled & 0.65 & 0.866049 & 0.436788 & 42.6 & 5.681533
\end{tabular}
\end{center}

Forecast skill improved:

\begin{center}
\begin{tabular}{lr}
contrast & forecast-skill delta \\
\hline
linear minus reactive & 0.033517 \\
linear minus shuffled & 0.045425 \\
nonlinear minus reactive & 0.084033 \\
nonlinear minus nonlinear-shuffled & 0.095941
\end{tabular}
\end{center}

But primary full-accounting residual-state predictability contrasts did not
pass:

\begin{center}
\begin{tabular}{lrrl}
contrast & mean delta & positive rate & interpretation \\
\hline
linear minus reactive & 0.114949 & 0.5 & inside label-permutation interval \\
linear minus shuffled & -0.053977 & 0.5 & inside label-permutation interval \\
nonlinear minus reactive & 0.077593 & 0.7 & inside label-permutation interval \\
nonlinear minus nonlinear-shuffled & -0.091333 & 0.3 & inside label-permutation interval
\end{tabular}
\end{center}

\paragraph{Promotion audit.}
The linear predictor passed forecast-skill and practical guardrails, but failed
residual-structure and oracle-nontriviality criteria. The nonlinear predictor
passed forecast-skill criteria, but failed practical guardrails and residual
structure criteria.

\paragraph{Decision.}
A5 was closed as a forecast-skill manipulation that did not demonstrate
promotion-worthy residual structured dynamics. The positive method result is
that prediction budget and forecast timing can be manipulated under matched
demand streams. The negative scientific result is that forecast improvement
alone did not produce residual structure beyond load, service, action
opportunity, task volume, and work-budget accounting.

\section{Cross-Cutting Lessons}

\subsection{Load and Service Variables Are Powerful Confounds}

Across A2--A5, many apparent structural effects tracked task inflow, queue
depth, service capacity, work opportunity, completion fraction, and action mix.
Future experiments must avoid endpoints that are just renamed forms of those
variables.

\subsection{Lobe Labels Are Diagnostics, Not Primary Evidence}

The existing lobe labels are built from queue and action fields. They can show
trajectory regimes, dwell, and transitions, but they cannot by themselves prove
independent collective structure. Future lobe-like observables should be less
directly tied to manipulated queue and action counts.

\subsection{Delay Matters, But Delay Alone Was Not Enough}

A4 delayed coupling produced a real raw completion-fraction synchrony
diagnostic. Once opportunity/load/action/transfer timing controls were applied,
the residual was no longer outside the null. Future delay experiments need
richer causal structure than transfer delay alone.

\subsection{Prediction Skill Is Not Emergent Dynamics}

A5 confirmed that bounded predictors can improve forecast skill under matched
demand streams. Residual structured dynamics did not appear under the current
scaffold. A future predictive experiment should make prediction a scarce
resource that competes with other activity, not just a controller input.

\subsection{The Missing Ingredient Is Nonlinear Mutual Dependence}

The experiments so far mostly varied pressure, service, scheduling, transfer
timing, and prediction budget. They did not yet implement strong nonlinear
mutual dependence among agents or hives. Logistic-shaped dependence of one
agent's activity on another directly targets this missing ingredient.

\section{Implications for Logistic-Dependence Experiments}

\subsection{Core Mechanism}

A next-stage model could make the activity of agent or hive $i$ depend
logistically on lagged activity of agent or hive $j$:

\[
a_i(t+1) =
\sigma\left(
  b_i
  + k_{ij}(x_j(t-\tau_{ij}) - \theta_{ij})
  + m_i(t)
  - c_i p_i(t)
\right)
\]

where:

\begin{itemize}
\item $a_i(t+1)$ is an activity probability, allocation share, or action bias.
\item $x_j(t-\tau_{ij})$ is lagged activity, load, prediction expenditure, or
state-burst intensity of another agent or hive.
\item $\sigma(z)=1/(1+\exp(-z))$ is the logistic nonlinearity.
\item $k_{ij}$ is the coupling slope.
\item $\theta_{ij}$ is a threshold near the typical operating range.
\item $\tau_{ij}$ is a delay.
\item $m_i(t)$ is memory or hysteresis.
\item $p_i(t)$ is prediction or modeling expenditure.
\item $c_i$ is the opportunity cost of prediction.
\end{itemize}

The most interesting regime is probably not saturation, where $a_i$ is near 0
or 1. It is the near-threshold regime where small changes in another agent's
lagged activity cause large but bounded changes in response.

\subsection{Why Logistic Dependence Is Different From Prior Knobs}

Prior knobs mostly changed exogenous pressure, service capacity, transfer
timing, or predictive control. Logistic dependence would introduce:

\begin{itemize}
\item threshold effects;
\item saturation;
\item nonlinear gain;
\item sensitivity to delays;
\item possible overcorrection;
\item possible hysteretic switching if memory is included;
\item bounded but sharp mutual influence among agents or hives.
\end{itemize}

Those ingredients are more plausible sources of structured recurrent dynamics
than simply increasing task pressure.

\subsection{Concrete Single-Hive Experiment Family}

Before jumping to three hives, a single hive could contain subpopulations or
roles whose activity depends logistically on each other.

\paragraph{State variables.}
For each role or task class, define coarse activity variables:

\begin{verbatim}
external_work_activity
research_activity
internal_improvement_activity
housekeeping_activity
prediction_activity
coordination_activity
\end{verbatim}

\paragraph{Example logistic couplings.}

\begin{itemize}
\item High external-work backlog increases housekeeping or coordination only
after a threshold.
\item Research starvation increases internal-improvement after a delay.
\item Too much internal-improvement suppresses external work after a threshold.
\item Prediction expenditure increases when other roles become hard to
forecast, but prediction expenditure reduces available work budget.
\end{itemize}

\paragraph{Suggested conditions.}

\begin{enumerate}
\item Reactive baseline with no logistic cross-role dependence.
\item Linear cross-role dependence with matched marginal gain.
\item Logistic same-tick dependence.
\item Delayed logistic dependence.
\item Delayed logistic dependence with hysteresis.
\item Delayed logistic dependence with bounded prediction-cost feedback.
\end{enumerate}

\paragraph{Primary comparison.}
The key contrast should not be throughput. It should be residual recurrent
structure after accounting for load, service, action opportunity, and total
work budget.

\subsection{Concrete Multi-Hive Experiment Family}

A three-hive ring is more suitable than a two-hive setup for meaningful target
and phase nulls. With two hives, direct and shuffled target assignment can be
equivalent because each hive has only one legal non-source target.

\paragraph{Minimal three-hive ring.}

\begin{verbatim}
hive A predicts and services pressure from hive B
hive B predicts and services pressure from hive C
hive C predicts and services pressure from hive A
\end{verbatim}

Each hive has bounded prediction budget and logistic response to lagged
activity from its neighbor:

\[
a_A(t+1)=\sigma(b_A+k_A(x_B(t-\tau)-\theta_A)-c_Ap_A(t))
\]

and similarly for B and C.

\paragraph{Useful variants.}

\begin{enumerate}
\item No coupling.
\item Linear ring coupling.
\item Logistic ring coupling without delay.
\item Logistic ring coupling with homogeneous delay.
\item Logistic ring coupling with heterogeneous delays.
\item Logistic coupling with prediction-cost feedback.
\item Logistic coupling with adaptive thresholds.
\end{enumerate}

\paragraph{Important controls.}

\begin{itemize}
\item Match transfer volume and source opportunity.
\item Use phase-shuffled coupling preserving marginal activity levels.
\item Use amplitude-matched linear controls.
\item Use threshold-shuffled controls preserving each hive's marginal
activation rate.
\item Use paired seeds and fixed exogenous demand streams.
\item Include no-coupling and delayed-but-nonpredictive controls.
\end{itemize}

\subsection{Candidate Observables}

Future observables should be chosen before implementation. Good candidates:

\begin{itemize}
\item Delay-embedded state vectors of role or hive activity after residualizing
load and service variables.
\item Recurrence plots over residual activity vectors.
\item Return-map structure in threshold-crossing intervals.
\item Motif counts over coarse macro-states not defined directly from queue
depth.
\item Divergence and reconvergence under tiny paired perturbations.
\item Nonlinear forecastability of macro-state sequences versus linear
forecastability.
\item Compression of macro-state strings relative to shuffled and
Markov-preserving nulls.
\item Phase relations among hives at preregistered lags.
\item Prediction expenditure bursts and their relation to later error,
overload, and recovery.
\end{itemize}

\subsection{Suggested Promotion Rule}

A logistic-dependence experiment should promote only if all of the following
hold:

\begin{enumerate}
\item Logistic coupling creates more residual recurrent structure than
amplitude-matched linear coupling.
\item The effect survives load, service, action, transfer-volume, and work
budget controls.
\item The effect survives phase-shuffled or threshold-shuffled nulls.
\item The effect is not merely worse queueing or longer backlog dwell.
\item Intermediate coupling slope or prediction budget is richer than both
zero coupling and saturated/oracle-like coupling.
\item Macro-state structure is partly predictable or compressible without being
trivially periodic.
\end{enumerate}

\section{Questions for a Strategic Reviewer}

\begin{enumerate}
\item Given the negative A2--A5 results, what is the smallest nonlinear
mechanism that would be meaningfully different from load/service accounting?
\item Should logistic dependence act on task creation, work allocation,
prediction expenditure, transfer probability, or role activation?
\item What observables are least confounded with queue depth and action count?
\item What is the right null for logistic threshold effects: phase shuffle,
threshold shuffle, amplitude-matched linear response, or Markov-preserving
macro-state shuffle?
\item Is a single-hive role-coupling model enough, or is a three-hive ring
needed for nontrivial phase and target nulls?
\item How should one distinguish structured strange-attractor-like recurrence
from a noisy limit cycle or from queue saturation?
\item What parameter ranges should be swept first: slope $k$, threshold
$\theta$, delay $\tau$, memory strength, prediction cost, or service capacity?
\item Can prediction budget be made an endogenous resource choice rather than
a fixed condition?
\end{enumerate}

\section{Bottom Line}

OmegaSim has so far taught a mostly negative but useful lesson. Pressure,
service capacity, scheduling, transfer timing, and forecast skill all produce
measurable effects, but apparent structure repeatedly collapses under
accounting controls. The next serious attempt at complex dynamics should avoid
another pressure or seed sweep. It should introduce a preregistered nonlinear
mutual-dependence mechanism, with logistic threshold response, delays, bounded
prediction cost, heterogeneity, and strong nulls.

The most promising next direction is:

\begin{quote}
Design a bounded single-hive or three-hive experiment in which role or hive
activity has logistic-shaped, delayed dependence on other role or hive
activity, with prediction effort as a costly managed resource, and test whether
intermediate slopes/budgets produce residual recurrent macro-structure beyond
load, service, action, and transfer accounting.
\end{quote}

\end{document}
